\hdhead{45}{多边形的认识——信息窗3}
\hdsection{1}{132}{填一填。}
\begin{scope}[scale=.6666667,draw=black,text=black]\draw[hdoutline] (382.000,308.000)--(419.000,203.000)--(513.000,203.000)--(550.000,308.000)--cycle;\end{scope}
\begin{scope}[scale=.6666667,draw=black,text=black]\draw[hdoutline,dashed] (419.000,203.000)--(419.000,308.000);\end{scope}
\begin{scope}[scale=.6666667,draw=black,text=black]\draw[line width=.45pt] (419.000,298.000)--(429.000,298.000)--(429.000,308.000);\end{scope}
\q{423}{178}{28}{13}{（}
\q{480}{178}{25}{13}{）}
\q{332}{249}{28}{13}{（}
\q{384}{249}{25}{13}{）}
\q{418}{249}{28}{13}{（}
\q{465}{249}{25}{13}{）}
\q{529}{249}{28}{13}{（}
\q{587}{249}{25}{13}{）}
\q{420}{315}{28}{13}{（}
\q{482}{315}{25}{13}{）}
\hdsection{2}{369}{分别画出下面梯形底边上的高。}
\begin{scope}[scale=.6666667,draw=black,text=black]\draw[hdoutline] (124.000,498.000)--(168.000,428.000)--(261.000,428.000)--(275.000,498.000)--cycle;\end{scope}
\begin{scope}[scale=.6666667,draw=black,text=black]\draw[hdoutline] (356.000,413.000)--(424.000,458.000)--(424.000,520.000)--(356.000,520.000)--cycle;\end{scope}
\begin{scope}[scale=.6666667,draw=black,text=black]\draw[hdoutline] (507.000,447.000)--(634.000,447.000)--(608.000,517.000)--(534.000,517.000)--cycle;\end{scope}
\begin{scope}[scale=.6666667,draw=black,text=black]\draw[hdoutline] (734.000,440.000)--(803.000,440.000)--(803.000,517.000)--(734.000,492.000)--cycle;\end{scope}
\begin{scope}[scale=.6666667,draw=black,text=black]\node[inner sep=0pt,font=\normalsize,text=black] at (201.000,516.000) {底};\end{scope}
\begin{scope}[scale=.6666667,draw=black,text=black]\node[inner sep=0pt,font=\normalsize,text=black] at (344.000,471.000) {底};\end{scope}
\begin{scope}[scale=.6666667,draw=black,text=black]\node[inner sep=0pt,font=\normalsize,text=black] at (571.000,431.000) {底};\end{scope}
\begin{scope}[scale=.6666667,draw=black,text=black]\node[inner sep=0pt,font=\normalsize,text=black] at (722.000,468.000) {底};\end{scope}
\hdsection{3}{559}{画一画，分一分。}
\q{102}{607}{730}{12}{① 在图形中各画一条线段，把图形分成一个平行四边形和一个三角形。}
\begin{scope}[scale=.6666667,draw=black,text=black]\draw[hdoutline] (300.000,715.000)--(322.000,647.000)--(399.000,647.000)--(440.000,715.000)--cycle;\end{scope}
\begin{scope}[scale=.6666667,draw=black,text=black]\draw[hdoutline] (559.000,648.000)--(685.000,648.000)--(642.000,715.000)--(583.000,715.000)--cycle;\end{scope}
\q{103}{738}{734}{12}{② 在每个图形中各画一条线段，把图形分成两个完全相同的梯形。}
\begin{scope}[scale=.6666667,draw=black,text=black]\draw[hdoutline] (301.500,846.000)--(339.500,776.000)--(447.500,776.000)--(409.500,846.000)--cycle;\end{scope}
\begin{scope}[scale=.6666667,draw=black,text=black]\draw[hdoutline] (571.000,776.000)--(680.000,776.000)--(680.000,846.000)--(571.000,846.000)--cycle;\end{scope}
\hdsection{4}{879}{在下面的点子图中分别画一个长方形、一个正方形、一个平行四边形和一个等腰梯形。}
\begin{scope}[scale=.6666667,draw=black,text=black]\foreach \i in {0,...,38}{\foreach \j in {0,...,9}{\fill[black] ({113.6+\i*18.4},{974.2+\j*18.4}) circle (1.15);}}\end{scope}
